\documentclass[12pt,a4paper]{article}

\usepackage{kotex}
\usepackage{amsmath,amssymb}
\usepackage{booktabs}
\usepackage{graphicx}
\usepackage[hidelinks,colorlinks=false]{hyperref}
\usepackage{array}
\usepackage{multirow}
\usepackage{geometry}
\usepackage{setspace}
\usepackage{caption}
\usepackage{subcaption}
\usepackage{xcolor}
\usepackage{tikz}
\usetikzlibrary{arrows.meta,shapes.geometric,positioning,fit,calc}
\usepackage{enumerate}
\usepackage{parskip}
\usepackage{tabularx}

\geometry{top=2.5cm, bottom=2.5cm, left=3cm, right=3cm}
\setstretch{1.4}
\captionsetup{font=small,labelfont=bf}
\captionsetup[table]{skip=4pt}
\renewcommand{\arraystretch}{1.12}
\setlength{\tabcolsep}{5pt}
\setlength{\emergencystretch}{2em}
\graphicspath{%
  {figures/new/ko/}%
  {figures/new/main/}%
  {figures/new/}%
  {figures/report/}%
  {figures/legacy/}%
}

% ─── 제목/저자 ────────────────────────────────────────────────────────
\title{%
  \textbf{TA-BRPL: 저전력 손실 네트워크에서\\
  Sinkhole 경로 장악 완화를 위한 신뢰 기반 BRPL}%
}
\author{이건형}
\date{2026년 3월}

\begin{document}
\maketitle
\thispagestyle{empty}

% ═══════════════════════════════════════════════════════════════════
\begin{abstract}
\noindent
\textbf{[문제]}
RPL 기반 저전력 손실 네트워크(LLN)에서 sinkhole 공격은
두 단계로 진행된다:
\emph{경로 장악 단계}---낮은 rank 위조로 트래픽을 조용히 끌어모아
네트워크 라우팅 구조를 지배하는---와
\emph{가용성 파괴 단계}---이후 선택적 패킷 드롭으로 PDR을 붕괴시키는.
기존 방어는 PDR 저하를 기준으로 탐지해 경로 장악 단계를 놓치거나,
신뢰 기반 즉각 차단으로 경로 교체(churn)가 폭증하는 문제를 안고 있다.

\textbf{[제안]}
본 연구는 Backpressure RPL(BRPL)의 부모 선택 점수에
경량 신뢰 패널티를 직접 결합한 \textbf{TA-BRPL}을 제안한다.
전달 행동($T_{fwd}$), 제어 평면 이상($T_{ctrl}$), 큐 정직성($T_{hon}$)
세 신호를 가중 기하 평균으로 집계하고,
비대칭 EWMA와 2단계 리뷰/블랙리스트·전환 게이트로
과도한 churn을 구조적으로 제어한다.

\textbf{[검증]}
Contiki-NG/Cooja에서 50노드 랜덤 토폴로지 75개 × 5 시드의
쌍체 실험(BRPL vs TA-BRPL)을 수행했다.
평가 1차 목표는 경로 장악 지표(att\_share/hit\_ratio), 2차 목표는
PDR 비열세, 3차 목표는 churn 상한 내 제어다.

\textbf{[결과]}
TA-BRPL은 공격자 경유율(att\_share)을 평균 15.4\% 감소
($\Delta = -0.0169$, Wilcoxon + Holm 보정 $p < 10^{-4}$,
$r_{rb} = 0.65$, 승률 76\%)시키고,
공격 구간 PDR을 1.0\,pp 개선했다 ($p = 0.002$).
churn 증가는 노드당 0.036회/공격구간이며,
부트스트랩 95\,\% CI 상한 $+0.057$이 사전 선언 상한 $+0.1$을 하회해
상한 내 비열세(H3)를 충족한다.

\textbf{[의의]}
본 연구는 방어의 1차 목표를 PDR이 아닌 경로 장악 억제로 재정의하고,
신뢰 기반 라우팅의 churn 폭증 문제를 안정화 설계 원칙으로
해결하는 체계를 제시한다.
다만 본 결론은 단일 공격자·UDGM·고정 공격 스케줄 조건에서의
근거이며, 단계 분리(stage-split)·오버헤드·오탐 지표는 확장 평가로 남겨 둔다.
\end{abstract}

\noindent\textbf{핵심어:}
LLN, RPL, BRPL, sinkhole, 신뢰 기반 라우팅, 경로 장악,
선택적 포워딩, Contiki-NG, Cooja

\clearpage
\tableofcontents
\clearpage


% ═══════════════════════════════════════════════════════════════════
\section{서론}
\label{sec:intro}

\subsection{문제: 조용히 시작하는 공격}

LLN에서 RPL은 센서 데이터의 수집·집계 품질을 좌우한다.
수십~수백 개의 배터리 기반 센서가
루트까지 최단·최저비용 경로를 찾아 데이터를 올려 보내는데,
이 경로 선택의 핵심은 ``rank가 낮은(루트에 가까운) 노드를 부모로 택한다''는
단순한 원칙이다.

\textbf{Sinkhole 공격}은 이 원칙을 정확히 겨냥한다.
공격자는 자신이 루트에 아주 가까운 것처럼 rank를 허위 광고하고,
주변 노드들이 자신을 부모로 채택하도록 유도한다.
여기서 결정적인 함정이 있다:
\emph{공격자는 초기에는 패킷을 정상적으로 전달한다.}
PDR은 멀쩡하다.
하지만 그 사이 점점 더 많은 노드가 공격자를 경유하도록 라우팅 구조가 재편된다.
이렇게 경로 장악(route capture)이 완성되면,
공격자가 언제든 선택적 드롭(grayhole/blackhole)을 시작할 수 있는
``준비된 위치''가 만들어진다.
이후 실제 드롭이 시작될 때는 이미 집중된 트래픽 구조 탓에
피해가 훨씬 크게 증폭된다~\cite{karlof2003,secrpl}.

\begin{figure}[t]
\centering
\includegraphics[width=0.95\linewidth]{fig_concept.pdf}
\caption{공격 단계 개념도(코드 스냅샷 기준).
  \textbf{Pre-attack}에서는 정상 합류가 이뤄지고,
  \textbf{Route capture}에서 낮은 rank DIO 유인으로 att\_share가 먼저 상승한다.
  이후 \textbf{Selective drop}이 활성화되면 PDR 저하가 뒤따른다.
  TA-BRPL의 1차 목표는 PDR 탐지 이전에 route capture를 억제하는 것이다.
  본 그림은 대표 동작을 설명하는 \emph{개념도}이며,
  메인 통계 결과 자체를 직접 시각화한 그래프는 아니다.}
\label{fig:concept}
\end{figure}

\textbf{문제의 핵심은 PDR로는 1단계를 볼 수 없다는 것이다.}
PDR 중심 탐지·평가 체계는 경로 장악이 완성된 뒤에야 이상을 감지한다.
이미 구조가 굳어진 상태에서 대응을 시작하면 늦다.

\subsection{기존 방어의 두 가지 한계}

기존 방어는 두 계열로 나뉜다.

\textbf{인증 기반 방식}(예: VeRA~\cite{dvir2011})은 서명된 DIO로 rank 조작을 차단한다.
그러나 공개키 인프라와 서명 검증 비용은 배터리로 구동되는
제한된 센서에는 과부하이고~\cite{roll_trust},
운영 중 키 배포/갱신 문제도 남는다.

\textbf{신뢰/이상 탐지 기반 방식}(예: SVELTE~\cite{raza2013})은 런타임에
악성 행동을 감지하려 하지만,
(a) 중앙 집중식 모니터에 의존하거나~\cite{iotroute_survey2016},
(b) 탐지 신호를 즉각적인 차단에 연결할 때
churn(과도한 경로 교체)이 폭증하는 부작용이 공통적으로 나타난다.
초기 버전 TA-BRPL에서도 이 문제가 관찰되었다(v1: churn +2.43).

\subsection{연구 질문과 가설}
\label{sec:rq}

위 배경으로부터 다음 두 연구 질문을 설정한다.

\begin{description}
  \item[RQ1 (Early-stage security)]
    sinkhole 공격의 경로 장악 단계에서,
    PDR 저하가 발생하기 \emph{전에}
    공격자 경유(att\_share/hit\_ratio)를 감소시킬 수 있는가?

  \item[RQ2 (Security–stability trade-off)]
    신뢰 신호를 라우팅 metric에 직접 결합할 때 발생하는
    churn 폭증을, 안정화 정책으로 통제하면서도
    공격 완화 효과를 유지할 수 있는가?
\end{description}

\noindent\textbf{RQ3(강건성/일반화)는 본문 핵심 검증 질문이 아니다.}
드롭률/온오프/다중 공격자/현실 무선모델에 대한 일반화는
본 라운드에서 ``확장 평가(extended evaluation)''로 분리해 기술한다.

이로부터 다음 세 가설을 사전 선언한다:

\begin{description}
  \item[H1 (격리 개선)]
    TA-BRPL은 BRPL 대비 공격 구간에서
    att\_share와 hit\_ratio를 감소시킨다.

  \item[H2 (PDR 비열세)]
    TA-BRPL의 공격 구간 PDR은 BRPL 대비 비열세(non-inferior)이다.
    비열세 마진: $\delta = -0.02$\,pp (사전 선언).

  \item[H3 (안정성 비용 제한 — 상한 내 비열세)]
    $\Delta$churn의 부트스트랩 95\,\% CI 상한이
    사전 선언한 실무적 허용 상한
    $\delta_{churn} = +0.1$\,회/노드/공격구간 미만이다
    (H2와 동일한 비열세 프레임으로 단방향 상한을 검증).
\end{description}

\noindent\textbf{지표 해석(중요).}
att\_share는 ``공격자 경유 트래픽 비율''로서 \emph{의존 강도}를,
hit\_ratio는 ``공격자를 한 번이라도 경유한 소스 비율''로서
\emph{오염 범위}를 측정한다.
두 지표를 함께 보고할 때 ``얼마나 많이 오염되었는가''와
``얼마나 넓게 퍼졌는가''를 분리해 해석할 수 있다.

\subsection{기여}

\begin{description}
  \item[C1 (문제 재정의)]
    sinkhole 방어의 1차 목표를 PDR이 아닌
    \emph{경로 장악 억제}(att\_share, hit\_ratio)로 재정의하고,
    이를 기본 지표로 채택하는 2단계 평가 틀을 제시한다.

  \item[C2 (경량 신뢰 통합)]
    전달($T_{fwd}$)·제어 평면($T_{ctrl}$)·큐 정직성($T_{hon}$)
    세 신호를 BRPL 부모 선택 점수에 소프트 패널티로 직접 통합한다.
    우리가 아는 한(to the best of our knowledge),
    별도 탐지 계층 없이 신뢰를 BRPL metric에 직접 결합한
    라우팅 내장형 설계의 \emph{초기 사례 중 하나}로 볼 수 있다.

  \item[C3 (안정화 설계 원칙)]
    비대칭 EWMA + 2단계 리뷰/블랙리스트 + 전환 마진/체류 게이트를
    ``신뢰 기반 라우팅의 churn 폭증을 구조적으로 제어하는 원칙''으로 제시한다.
    단순 파라미터 튜닝이 아니라 설계 원칙의 일반화이다.

  \item[C4 (분산 평가 방법론)]
    단일 토폴로지 편향을 피하기 위해
    밀도별 층화 75개 랜덤 토폴로지와 쌍체 통계(Wilcoxon + Holm, 부트스트랩 CI)를
    적용해 효과의 \emph{일관성}을 검증한다.
\end{description}


% ═══════════════════════════════════════════════════════════════════
\section{배경 및 관련 연구}
\label{sec:background}

\subsection{표준 축: RPL과 BRPL}

RPL~\cite{rfc6550}은 주기적 DIO 메시지로 DODAG를 구성하며,
표준 목적함수로 MRHOF~\cite{rfc6719}, OF0~\cite{rfc6552}가 정의된다.
본 연구는 MRHOF(ETX 기반)와 BRPL을 비교축으로 두고,
메트릭 컨테이너/경로 비용 표현은 RFC~6551~\cite{rfc6551}의 틀을 따른다.
BRPL~\cite{brpl}은 RPL에 역압(backpressure) 항을 추가한다:
\begin{equation}
  S_{\text{BRPL}}(v) = w_p \cdot \text{rank}(v)
                     + w_q \cdot \Delta Q(v)
                     + w_l \cdot \text{ETX}(v).
  \label{eq:brpl}
\end{equation}
$\Delta Q(v)$는 로컬 노드와 $v$의 큐 점유 기울기로,
혼잡 노드를 자연스럽게 우회하는 복원력을 제공한다.
그러나 공격자가 큐 광고를 허위로 낮추면 이 방어가 무력화된다
---이것이 TA-BRPL의 $T_{hon}$ 신호가 채우는 공백이다.

\subsection{공격 축: Sinkhole + 선택적 포워딩}

IETF의 RPL 보안 위협 분석(RFC~7416~\cite{rfc7416})은
sinkhole을 ``topology manipulation''으로 공식 분류하며,
rank 위조를 통한 DODAG 경로 조작을 핵심 위협으로 지정한다.
Karlof and Wagner~\cite{karlof2003}는 sinkhole 공격이
유인(luring) → 장악(capture) → 드롭(drop)의 순서로 진행됨을 보였다.
최근 분류 연구~\cite{secrpl}는 이를 ``topology manipulation''과
``selective forwarding''의 복합 공격으로 체계화한다.
공통적인 평가 공백으로 지적되는 것이 바로
``PDR만 측정해 장악 단계를 놓치는 문제''다.
Marti 등~\cite{marti2000}은 grayhole에서 ACK 정상 처리가
단순 모니터링으로는 탐지를 어렵게 한다는 점을 밝혔다.

\medskip
\noindent\textbf{코드 스냅샷 기준 공격 모델.}
본 원고의 공격 시나리오는
\path{motes/sinkhole_drop_attacker.c} 구현과
\path{Makefile.sinkhole_drop_*} 빌드 플래그를 기준으로 고정했다.
공격 노드는 먼저 정상 DODAG에 합류하고,
$t \ge 350$\,s (\texttt{ATTACK\_WARMUP\_SECONDS})부터
(i) 15\,s 주기(\texttt{SINKHOLE\_DIO\_PERIOD\_SECONDS})로
\texttt{ROOT\_RANK + 1} (\texttt{SINKHOLE\_RANK\_DELTA=1})의 저-rank DIO를 송신하며,
(ii) 루트향 포워딩 UDP(포트 8765)에 대해
$p=\texttt{ATTACK\_DROP\_PCT}/100$ 확률로 선택적 드롭을 수행한다.
즉 드롭률은 0/25/50/75/100\% 스윕이 가능하고,
제어 패킷(DIO/DAO/DIS)은 드롭 대상에서 제외된다.

\subsection{방어 축: 세 계열과 TA-BRPL의 위치}

\textbf{(a) 암호/인증 기반}: VeRA~\cite{dvir2011} 등 서명 기반 DIO 보호.
리소스 오버헤드와 인프라 의존성이 제한된 LLN에의 적용을 어렵게 한다.

\textbf{(b) IDS/중앙 탐지 기반}: SVELTE~\cite{raza2013} 등
런타임 이상 탐지. 중앙 모니터 또는 대역 외 채널이 필요하며,
라우팅 결정 자체를 바꾸지는 못한다.

\textbf{(c) 신뢰/평판 기반 라우팅 내장형}:
Ganeriwal 등~\cite{airtime_trust}의 watchdog 평판과
Momani 등~\cite{Momani2010}의 WSN 신뢰 모델이 기초를 놓았다.
최근 RFTrust(2021)~\cite{rftrust2021},
PCC-RPL(2022)~\cite{pccrpl2022},
RPL 신뢰 라우팅 모델(Sensors 2022)~\cite{trustmodel2022}은
신뢰 기반 대응의 실용성을 확장했다.
그러나 다수 연구는 신뢰 추정 품질 중심으로 보고하며,
``신뢰-의사결정 결합 시 churn을 어떻게 정책적으로 제어하는가''는
상대적으로 덜 다뤄진다~\cite{sinkhole_survey2023,rpl_survey2025}.
ROLL 워킹그룹의 신뢰 아키텍처~\cite{roll_trust}는
분산 신뢰 통합의 방향성을 제시하나 아직 표준화 전이다.

\textbf{TA-BRPL의 차별점}: (i)~신뢰를 BRPL metric에 \emph{직접 통합}하여
탐지/라우팅 분리를 줄이고, (ii)~전달·제어 평면·큐 정직성 세 신호를
동시에 반영하며, (iii)~admission/retention/re-entry 게이트로
decision-level stabilization을 명시적으로 설계한다.


% ═══════════════════════════════════════════════════════════════════
\section{TA-BRPL 설계}
\label{sec:design}

\subsection{핵심 원칙: 소프트 패널티 + 보수적 게이트}

TA-BRPL의 설계 철학은 두 문장으로 요약된다:
\begin{quote}
\textit{``의심스러운 노드는 즉시 차단하지 말고, 조용히 덜 선호하라.''}\\
\textit{``행동을 바꿀 때는 충분한 증거가 쌓인 뒤에 천천히 바꿔라.''}
\end{quote}
첫 문장은 소프트 패널티 통합 방식을 정당화하고,
두 번째 문장은 안정화 게이트 설계의 근거가 된다.

\begin{figure}[t]
\centering
\begin{tikzpicture}[
  node distance=0.55cm and 1.6cm,
  box/.style={rectangle,draw,fill=blue!8,rounded corners=3pt,
              minimum width=36mm,minimum height=9mm,align=center,font=\small},
  sbox/.style={rectangle,draw,fill=orange!12,rounded corners=3pt,
               minimum width=36mm,minimum height=9mm,align=center,font=\small},
  lbox/.style={rectangle,draw,fill=orange!12,rounded corners=3pt,
               minimum width=55mm,minimum height=14mm,align=center,font=\scriptsize},
  agg/.style={rectangle,draw,fill=green!12,rounded corners=3pt,
              minimum width=36mm,minimum height=11mm,align=center,font=\small},
  arrow/.style={-{Stealth},thick},
]
% 왼쪽 열: 세 신뢰 신호
\node[box] (tf)                     {$T_{fwd}$: 전달 신뢰도};
\node[box, below=of tf] (tc)        {$T_{ctrl}$: 제어 평면 신뢰};
\node[box, below=of tc] (th)        {$T_{hon}$: 큐 정직성};
% 가운데 열: 집계 + 위아래 정책
\node[agg, right=of tc]  (ta)       {$T_{agg}$\\가중 기하 평균};
\node[sbox, above=of ta] (ew)       {비대칭 EWMA\\(빠른 하락, 느린 회복)};
\node[sbox, below=of ta] (gt)       {안정화 게이트\\마진·체류·$\tau_{join}$};
% 오른쪽 열: BRPL 점수
\node[lbox, right=of ta] (br)       {BRPL 점수 반영(식 \ref{eq:trust_penalty}, \ref{eq:trust_scales})\\
  $S'_{TA}=S_{BRPL}\cdot\frac{t}{1+\lambda_p(1-t)}$\\
  $S_{TA}=S'_{TA}\cdot\eta_{val}\cdot\eta_{policy}\cdot\eta_{sticky}$};
% 화살표
\draw[arrow](tf.east)--++(0.4,0)|-(ta.west);
\draw[arrow](tc.east)--(ta.west);
\draw[arrow](th.east)--++(0.4,0)|-(ta.west);
\draw[arrow](ta.north)--(ew.south);
\draw[arrow](ta.south)--(gt.north);
\draw[arrow](ta.east)--(br.west);
\end{tikzpicture}
\caption{TA-BRPL 신뢰 파이프라인.
  세 신호가 집계 점수로 통합되고,
  비대칭 EWMA로 업데이트된 뒤 BRPL 점수에 소프트 패널티로 반영된다.
  안정화 게이트는 과잉 반응을 차단하는 마지막 방어선이다.}
\label{fig:arch}
\end{figure}

\subsection{세 가지 신뢰 신호}

\noindent\textbf{$T_{fwd}$ — 전달 신뢰도.}
에코 기반 수동 도청(passive overhearing)으로 측정한다.
로컬 노드가 $v$에게 패킷을 전송한 뒤 $v$의 재전송 소리를 들으면 성공,
들리지 않으면 실패로 기록한다:
\begin{equation}
  T_{fwd}(v) = \frac{\text{succ}(v) + \alpha}{\text{total}(v) + \alpha + \beta},
\end{equation}
$\alpha, \beta$는 약한 Beta 사전 분포 파라미터이다.
공격자가 패킷을 버리면 에코가 오지 않아 $T_{fwd}$가 하락한다.

\noindent\textbf{$T_{ctrl}$ — 제어 평면 신뢰.}
세 가지 DIO 이상을 선형 결합해 sinkhole 유인을 탐지한다:
불가능한 rank 주장($A_{rank}$),
지속적 저-rank 유인($A_{dio}$),
버전 번호 비일관성($A_{ver}$):
\begin{equation}
  T_{ctrl}(v)=1-\tfrac{w_r A_{rank}+w_d A_{dio}+w_v A_{ver}}{10}.
\end{equation}

\noindent\textbf{$T_{hon}$ — 큐 정직성.}
$v$가 광고하는 큐 깊이와 포워딩 지연으로 추정한 큐 깊이를 비교한다:
\begin{equation}
  T_{hon}(v)=1-\min\!\left(1,\;
    \frac{|Q_{adv}(v)-Q_{est}(v)|}{Q_{max}}\right).
\end{equation}
백로그가 있지만 빈 큐를 광고하는 공격자를 포착한다.
이 신호는 BRPL 기반 설계에서 특히 중요한데,
BRPL이 큐 정보를 신뢰한다는 설계 가정이 있기 때문이다.

\noindent\textbf{집계.} 세 신호를 가중 기하 평균으로 통합한다:
\begin{equation}
  T_{agg}(v,t) =
  T_{fwd}(v,t)^{w_f}\cdot
  T_{ctrl}(v,t)^{w_c}\cdot
  T_{hon}(v,t)^{w_h},
  \quad
  w_f+w_c+w_h=1.
  \label{eq:tagg}
\end{equation}
코드 스냅샷의 실장값은 $(w_f,w_c,w_h)=(0.5,0.3,0.2)$이다.
기하 평균을 쓰는 이유는 \textit{단일 신호 gaming}에 대한 강건성 때문이다:
한 성분만 비정상적으로 높아도 나머지 성분이 낮으면 전체 점수가
자동으로 제한되어 arithmetic mean 대비 과신을 줄인다.

\subsection{비대칭 EWMA와 2단계 검증}

신뢰는 60초마다 비대칭 EWMA로 업데이트된다:
\begin{equation}
  \hat{T}^{(t+1)}(v)=
  \begin{cases}
    \lambda_{down}\cdot \hat{T}^{(t)}(v)
      +(1-\lambda_{down})\cdot T_{agg}(v,t),
      &T_{agg}(v,t) < \hat{T}^{(t)}(v)\\
    \lambda_{up}\cdot \hat{T}^{(t)}(v)
      +(1-\lambda_{up})\cdot T_{agg}(v,t),
      &\text{그 외}\;\text{(회복)}
  \end{cases}
  \label{eq:ewma}
\end{equation}
$\lambda_{down} = 0.4 < \lambda_{up} = 0.7$으로,
악화는 빠르게 반영하되 회복에는 더 많은 증거를 요구한다.
일시적 링크 불량에 의한 거짓 양성을 줄이는 핵심 설계다.

\medskip
\noindent\textbf{BRPL 점수 반영(실장식).}
TA-BRPL은 신뢰를 별도 탐지 알람으로 쓰지 않고
BRPL 후보 비용에 직접 반영한다.
기본 비용 $S_{BRPL}(v)$에 대해 정규화 신뢰
$t(v)=\hat{T}(v)/1000$와 불신뢰도 $d(v)=1-t(v)$를 두면,
코드의 핵심 반영식은 다음과 같다(선형 $\gamma=1$ 기준):
\begin{equation}
  S'_{TA}(v)=
  S_{BRPL}(v)\cdot
  \frac{t(v)}{1+\lambda_p d(v)},
  \label{eq:trust_penalty}
\end{equation}
\begin{equation}
  S_{TA}(v)=
  S'_{TA}(v)\cdot
  \eta_{val}(v)\cdot
  \eta_{policy}(v)\cdot
  \eta_{sticky}(v).
  \label{eq:trust_scales}
\end{equation}
그림~\ref{fig:arch}의 우측 블록 수식은 식~\eqref{eq:trust_penalty},
\eqref{eq:trust_scales}와 동일한 반영식만을 요약 표기한 것이다.
여기서 $\eta_{val}$은 review/penalized 상태의 검증 스케일,
$\eta_{policy}$는 re-entry/escape/hold 기반 추가 스케일,
$\eta_{sticky}$는 불필요한 즉시 재전환을 막는 current-parent stickiness이다.
코드 스냅샷에서는 $\lambda_p = \texttt{BRPL\_CONF\_TRUST\_LAMBDA\_PENALTY}/1000 = 0.45$를 사용한다.
또한 hard blacklist는 후보 집합에서 직접 제외된다.

신뢰 수준에 따른 2단계 검증:
$T_{agg} < \tau_{warn}=0.60$이면 \emph{리뷰} 단계 진입,
3회 관측 창 중 2회 이상 불량이 확인될 때에만
120초 \emph{블랙리스트}로 격상된다.
하드 제외 전에 약 3분의 증거가 필요하도록 설계해
정직한 노드의 오탐 블랙리스팅을 억제한다.

\subsection{안정화 게이트: Churn 폭증 방지}

표~\ref{tab:versions}에서 보듯, 게이트 없는 초기 버전(v1)은
churn이 +2.43으로 폭증해 격리 이득을 상쇄했다.
세 가지 게이트가 이를 막는다:

\begin{enumerate}[(1)]
  \item \textbf{전환 마진 게이트}: 새 후보가 현재 부모보다
    상대 17\% \emph{이상 그리고} 절대 90단위 \emph{이상} 좋아야만 전환.
  \item \textbf{체류 시간 게이트}: 최소 180초간 현재 부모 유지.
  \item \textbf{허용 임계값}: 신규 후보는 $T_{agg}\ge\tau_{join}=0.51$,
    현재 부모는 $T_{agg}\ge\tau_{warn}=0.60$ 이하 시 패널티 가중치 증가.
\end{enumerate}

\begin{figure}[t]
\centering
\fbox{\parbox{0.94\linewidth}{
\textbf{Algorithm 1. TA-BRPL 부모 전환 게이트 (요약)}\\[2pt]
\texttt{Input: 현재 부모 $p$, 도전자 $c$, 비용 $S_{TA}$, 신뢰 상태}\\
\texttt{1) if $c$ is blacklisted or admission-fail: return $p$}\\
\texttt{2) if dwell-time not satisfied and no escape exception: return $p$}\\
\texttt{3) $\Delta \leftarrow S_{TA}(p)-S_{TA}(c)$}\\
\texttt{4) if $\Delta < \Delta_{abs}+\Delta_{extra}$ and $\Delta/S_{TA}(p) < \Delta_{rel}$: return $p$}\\
\texttt{5) else: return $c$}
}}
\caption{게이트 정책 의사코드. 핵심은 ``의심 후보 즉시 차단''이 아니라
  admission/dwell/margin을 순차 적용해 불필요한 재전환을 억제하는 것이다.}
\label{fig:algo_gate}
\end{figure}

\noindent\textbf{용어 정의(재현성).}
Algorithm 1에서 \emph{admission-fail}은
``후보 부모가 블랙리스트 상태이거나 신규 부모 허용 임계
($T_{agg}<\tau_{join}$)를 만족하지 못한 경우''를 뜻한다.
\emph{dwell-time}은 현재 부모 최소 유지시간(기본 180초)이고,
\emph{escape exception}은 현재 부모가 공격자 의심 상태로 누적될 때
보호 탈출 전환을 허용하는 예외 경로를 의미한다.
각 용어의 상태 변수/파라미터 고정값은 \S\ref{app:repro_lock},
로그 기반 집계 창 정의는 \S\ref{app:metrics}와 1:1로 연결한다.

\begin{table}[t]
\centering
\caption{설계 버전별 평균 성능 격차
  ($\Delta = \text{TA-BRPL}-\text{BRPL}$, 고정 토폴로지 GRID+BOTTLE 40회).
  굵은 줄이 메인 실험에 사용된 최종 버전이다.}
\label{tab:versions}
\begin{tabular}{lrrrr}
\toprule
버전(핵심 변경) & $\Delta$att & $\Delta$hit & $\Delta$churn & $\Delta$PDR \\
\midrule
v1 (게이트 없음)              & $+0.013$ & $+0.048$ & $+2.43$ & $-0.132$ \\
v3 (허용 제어 도입)           & $+0.004$ & $+0.032$ & $+0.69$ & $-0.034$ \\
v13.8 (손실 모드 게이트)      & $-0.042$ & $-0.015$ & $+0.47$ & $-0.001$ \\
\textbf{v13.12 (탈출 쿨다운)} & $\mathbf{-0.042}$ & $\mathbf{-0.015}$
                               & $\mathbf{+0.42}$ & $\mathbf{-0.001}$ \\
\bottomrule
\end{tabular}
\end{table}

\subsection{공격자 구현 파라미터(코드 스냅샷)}
\label{app:attack_impl}

\begin{table}[htbp]
\centering
\caption{\texttt{sinkhole\_drop\_attacker.c} 기반 공격 파라미터}
\label{tab:attack_impl}
\begin{tabular}{ll}
\toprule
항목 & 값 (코드 기준) \\
\midrule
워밍업 (\texttt{ATTACK\_WARMUP\_SECONDS}) & 350 s \\
가짜 DIO 주기 (\texttt{SINKHOLE\_DIO\_PERIOD\_SECONDS}) & 15 s \\
rank 위조 폭 (\texttt{SINKHOLE\_RANK\_DELTA}) & +1 (루트 대비) \\
드롭 대상 & 포워딩 UDP to root (port 8765) \\
드롭 확률 (\texttt{ATTACK\_DROP\_PCT}) & 0/25/50/75/100\% sweep 가능 \\
제어 패킷(DIO/DAO/DIS) & 드롭 제외 \\
\bottomrule
\end{tabular}
\end{table}


% ═══════════════════════════════════════════════════════════════════
\section{실험 결과}
\label{sec:results}

\subsection{한 줄 요약}

\begin{center}
\fbox{\parbox{0.88\linewidth}{
  H1 채택: att\_share/hit\_ratio 유의하게 감소 ($p < 10^{-4}$).\\
  H2 채택: PDR 비열세 93.3\% 통과, 집단 수준 CI 하한 $> -0.02$.\\
  H3 채택: churn 증가 +0.036/노드/구간, 사전 선언 상한 0.1 이내.
}}
\end{center}

\subsection{메인 실험 전 4-way 파이프라인 사전 검증}
\label{sec:precheck4way}

메인 75$\times$5 전수 실험에 앞서,
\texttt{run\_random\_topo\_drop\_precheck.sh}로
동일 파이프라인/동일 정책 스냅샷의 4-way(RPL/BRPL/SMTRUST/TABRPL)
재실행이 정상 동작함을 확인했다.
사전 검증의 목적은 \textbf{성능 우열 결론}이 아니라
\textbf{실험 파이프라인 무결성 확인}이다.
사전 스냅샷 수치(부분 완료 상태)는 부록 표~\ref{tab:precheck4way_appendix}에 제시한다.
이는 ``코드/시나리오/로그/집계'' 체인 점검용이며,
본 논문의 통계적 결론은 이 precheck가 아니라,
동일 조건 전수 완료 후의 본 실험 결과(이하 섹션)를 따른다.

\subsection{프로토콜별 절대값 비교}

\begin{table}[t]
\centering
\caption{프로토콜별 공격 구간 절대값 (랜덤 토폴로지 실험, 공격 구간 기준).
  주 비교 대상은 BRPL과 TA-BRPL이다.
  RPL은 기저 참조값으로 제시하며, SMTRUST는 \S\ref{sec:precheck4way}에서
  동일 파이프라인 사전 검증 결과로 별도 보고한다.}
\label{tab:absolute}
\begin{tabular}{lrrrr}
\toprule
프로토콜 & att\_share & hit\_ratio & PDR\_dur & churn \\
\midrule
RPL (MRHOF)     & 0.110 & 0.117 & 0.951 & 0.037 \\
BRPL            & 0.110 & 0.117 & 0.950 & 0.037 \\
\textbf{TA-BRPL}& \textbf{0.093} & \textbf{0.101} & \textbf{0.961} & 0.073 \\
\bottomrule
\end{tabular}
\end{table}

표~\ref{tab:absolute}는 주요 3개 프로토콜의 공격 구간 절대값이다.
RPL(MRHOF)과 BRPL의 att\_share/hit\_ratio가 거의 동일한 것은,
BRPL의 역압 메커니즘만으로는 sinkhole 경로 장악을 억제하기에
부족함을 보여준다.
TA-BRPL은 이를 15.4\% 개선한다.

\subsection{Baseline 비교축의 타당성}

본 연구의 baseline 구성은 단순 병렬 비교가 아니라
문헌 포지셔닝을 위한 3축 비교를 목표로 설계했다:
(i) 표준 라우팅축(RPL), (ii) 혼잡 인지축(BRPL),
(iii) 동급 보안축(RFTrust/PCC-RPL 계열).
현재 본문 통계 결론은 BRPL 중심 쌍체 비교로 닫고,
동급 보안 baseline은 본실험 확장(F1)에서 공정 비교로 다룬다.
SMTRUST는 본문 결론 baseline이 아니라
파이프라인 무결성 점검(\S\ref{sec:precheck4way}) 용도로만 사용한다.

\begin{table}[t]
\centering
\caption{Baseline 축별 역할과 현재 증거 수준}
\label{tab:baseline_axis_status}
\begin{tabular}{p{2.5cm}p{5.0cm}p{4.8cm}}
\toprule
축 & 역할 & 현재 상태 \\
\midrule
RPL (MRHOF) & 표준 ETX 기반 기저선 & 본문 절대값 비교에 포함 (표~\ref{tab:absolute}) \\
BRPL & 혼잡 인지 라우팅 기저선(주 비교축) & 75개 토폴로지 쌍체 통계로 결론 확정 (표~\ref{tab:stats}) \\
RFTrust/PCC-RPL 계열 & 동급 보안 라우팅 비교축 & 구현 정합성 확보 후 본실험 확장(F1)에서 직접 비교 예정 \\
SMTRUST (sanity only) & 파이프라인 스모크 테스트 축 & precheck/부록 한정 보고 (\S\ref{sec:precheck4way}) \\
\bottomrule
\end{tabular}
\end{table}

\subsection{신뢰 성분 Ablation (pre-main 완료)}
\label{sec:ablation_plan}

세 신호($T_{fwd},T_{ctrl},T_{hon}$)를 핵심 기여로 주장하기 위해,
최소 구성 ablation을 본문 체계에 포함한다.
메인 75쌍 실험 전, 동일 파이프라인/동일 정책 스냅샷에서
단축 pre-main 검증(3 density $\times$ 3 topology seed $\times$ 1 run seed,
총 9 topology-paired 비교)을 완료했다
(\texttt{results/random\_topo\_ablation\_minset\_v1}).

\begin{table}[t]
\centering
\caption{최소 ablation 구성(동일 정책 스냅샷 기준)}
\label{tab:ablation_minset}
\begin{tabular}{lccc}
\toprule
구성 & $w_f$ & $w_c$ & $w_h$ \\
\midrule
Full (기준) & 0.5 & 0.3 & 0.2 \\
$T_{fwd}$ only & 1.0 & 0.0 & 0.0 \\
$T_{fwd}{+}T_{ctrl}$ & 0.7 & 0.3 & 0.0 \\
\bottomrule
\end{tabular}
\end{table}

\begin{table}[t]
\centering
\caption{최소 ablation pre-main 결과(전체 9 topology-paired, $\Delta$ = variant $-$ BRPL)}
\label{tab:ablation_minset_results}
\begin{tabular}{lccccc}
\toprule
구성 & $\Delta$att\_share [95\% CI] & $\Delta$hit\_ratio & $\Delta$churn & $\Delta$PDR\_dur & Win(att) \\
\midrule
Full (기준) & $-0.0335$ [$-0.0473$, $-0.0186$] & $-0.0294$ & $+0.0850$ & $+0.0048$ & 88.9\% \\
$T_{fwd}$ only & $-0.0073$ [$-0.0263$, $+0.0114$] & $+0.0065$ & $+0.1438$ & $-0.0337$ & 55.6\% \\
$T_{fwd}{+}T_{ctrl}$ & $-0.0076$ [$-0.0274$, $+0.0114$] & $+0.0000$ & $+0.0948$ & $-0.0304$ & 55.6\% \\
\bottomrule
\end{tabular}
\end{table}

표~\ref{tab:ablation_minset_results}에서 Full 모델만
$\Delta$att\_share의 95\% CI가 0 아래로 완전히 분리된다.
반면 $T_{fwd}$ only와 $T_{fwd}{+}T_{ctrl}$는 CI가 0을 가로지르고,
PDR\_dur 평균도 각각 $-0.0337$, $-0.0304$로 하락한다.
즉 pre-main 수준에서도 ``세 신호 결합 + 안정화 정책''이
격리 이득과 전달 품질의 동시 확보에 필요함을 확인했다.

\subsection{전체 통계 결과 (H1·H2·H3 검증)}

\begin{table}[t]
\centering
\caption{전체 쌍체 통계 검정 결과 (75개 토폴로지 쌍, 부트스트랩 95\% CI).
  H1/H2: 단측 Wilcoxon + Holm-Bonferroni 보정; $^\dagger$: Holm 보정 후 $p < 0.05$.
  H3: CI 상한 비열세 검정; $^\ddagger$: CI 상한 $+0.057 < \delta_{churn} = +0.1$ 충족.}
\label{tab:stats}
\begin{tabular}{llrrl}
\toprule
가설 & 지표 & 평균 $\Delta$ & 95\% CI & $p$ (Holm) \\
\midrule
H1 & att\_share (↓) & $-0.0169$ & $[-0.023,\,-0.011]$ & $<10^{-4}$$^\dagger$ \\
H1 & hit\_ratio (↓) & $-0.0166$ & $[-0.024,\,-0.009]$ & $<10^{-4}$$^\dagger$ \\
H2 & PDR\_dur   (↑) & $+0.0102$ & $[+0.003,\,+0.018]$ & $0.002$$^\dagger$ \\
H3 & churn (상한$<$0.1) & $+0.0364$ & $[+0.018,\,+0.057]$ & CI 상한 채택$^\ddagger$ \\
\bottomrule
\end{tabular}
\end{table}

표~\ref{tab:stats}와 그림~\ref{fig:delta_ci}은 H1·H2·H3를 검증한다.

\subsection{메인 지표의 분포 증거}

\begin{figure}[t]
\centering
\includegraphics[width=0.95\linewidth]{fig_main_delta_box4.pdf}
\caption{메인 지표의 쌍체 $\Delta$ 분포(Boxplot, 75개 토폴로지).
  점선은 $\Delta=0$ 기준선이다.
  att\_share와 hit\_ratio는 분포 중심이 0 아래로 이동하고,
  PDR\_dur는 0 위쪽으로 이동한다.
  churn은 양의 방향으로 이동하지만 폭은 상대적으로 제한적이다.}
\label{fig:delta_box4}
\end{figure}

그림~\ref{fig:delta_box4}는 평균값뿐 아니라
토폴로지별 분포 전체가 이동했음을 보여준다.
즉 TA-BRPL의 이득은 특정 outlier 토폴로지의 효과가 아니라
쌍체 분포 수준의 일관된 이동으로 해석된다.

\medskip
\noindent\textbf{H1 (격리 개선) — 채택.}
att\_share와 hit\_ratio가 Holm 보정 후에도 강한 유의성을 보인다
($p < 10^{-4}$, $r_{rb} = 0.65$, 중-대형 효과 크기).
95\% CI 전체가 0 왼쪽에 위치해 방향 일관성을 확인한다.
여기서 att\_share 감소는 ``공격자 의존 강도'' 완화를,
hit\_ratio 감소는 ``오염 범위'' 축소를 의미한다.
다만 본 실험의 주 평가 창(350--650초)은 sinkhole 유인과 드롭이 동시 활성인 창이므로,
``드롭 이전 장악 단계''를 독립적으로 분리한 인과 검증은
확장 평가(stage-split)로 별도 수행한다.

\medskip
\noindent\textbf{H2 (PDR 비열세) — 채택.}
$\Delta\text{PDR} = +0.0102$이고 CI 하한 $+0.003 > -0.02$ (비열세 마진),
집단 수준에서 비열세가 확인된다.
토폴로지별로는 93.3\%가 통과(그림~\ref{fig:noninf}).

\medskip
\noindent\textbf{H3 (상한 내 비열세) — 채택.}
H3의 판정 기준은 ``$\Delta$churn의 95\,\% CI 상한 $< +0.1$''이다.
표~\ref{tab:stats}에서 CI 상한 $+0.057 < +0.1$이므로 H3가 충족된다.
절대 증가량 $+0.036$회/노드/구간은 사전 선언 상한의 36\% 수준으로,
격리 이득 대비 \emph{관리 가능한 안정성 오버헤드}
(controlled stability overhead)로 해석할 수 있다.
(참고: 동일 데이터를 양측 Wilcoxon으로 보면 $p \approx 0.02$로 소폭 증가가 통계적으로
유의하지만, 효과 크기는 매우 작고 실무 상한을 초과하지 않는다.)

\begin{figure}[t]
\centering
\includegraphics[width=0.92\linewidth]{fig4_delta_ci.pdf}
\caption{$\Delta$ (TA-BRPL $-$ BRPL)의 부트스트랩 95\% 신뢰구간 (75개 토폴로지 쌍).
  att\_share, hit\_ratio의 CI는 0 왼쪽(H1 채택),
  PDR\_dur의 CI는 0 오른쪽(H2 채택),
  churn의 CI는 0 우측에 걸쳐 소폭 증가(H3 상한 내 채택).}
\label{fig:delta_ci}
\end{figure}

\subsection{밀도별 세부 결과}

\begin{table}[t]
\centering
\caption{밀도별 결과 (각 25개 토폴로지 쌍).
  $^\dagger$: Holm 보정 $p < 0.05$.}
\label{tab:density}
\begin{tabular}{lrrrrl}
\toprule
밀도 & $\Delta$att [CI] & 승률 & $\Delta$PDR & $\Delta$churn & $p$(att) \\
\midrule
고밀도 & $-0.009$ [$-0.016,\,-0.001$] & 68\% & $+0.002$ & $-0.001$ & $0.206$ \\
중밀도 & $-0.026$ [$-0.038,\,-0.015$] & \textbf{84\%} & $+0.012$$^\dagger$ & $+0.021$ & $<10^{-3}$$^\dagger$ \\
저밀도 & $-0.016$ [$-0.027,\,-0.004$] & 76\% & $+0.016$ & $+0.089$ & $0.048$$^\dagger$ \\
\bottomrule
\end{tabular}
\end{table}

그림~\ref{fig:attshare_box}--\ref{fig:density_bar}은 밀도별 결과를 시각화한다.

\medskip\noindent\textbf{중밀도가 최적인 이유.}
중밀도에서 att\_share 감소 26\%(\textbf{승률 84\%})로 가장 강한 효과가 나타난다.
충분한 경로 다양성이 있어 신뢰 모델이 트래픽을 우회시킬 수 있되,
고밀도처럼 경로 중복성 자체가 공격자 영향을 희석시키지는 않는 구간이다.

\medskip\noindent\textbf{고밀도에서 유의성이 낮은 이유.}
고밀도에서는 토폴로지 구조 자체가 이미 공격자 영향을 분산시켜
신뢰 모델의 한계 이득이 줄어든다 (승률 68\%, Holm 보정 후 비유의).

\medskip\noindent\textbf{저밀도에서 churn 비용.}
경로 선택지가 적어 우회 재구성이 더 자주 일어나
churn 비용이 $+0.089$로 가장 높다.
그러나 att\_share 개선(16\%)은 유효하다.

\begin{figure}[t]
\centering
\includegraphics[width=0.92\linewidth]{fig3_att_share_box_by_density.pdf}
\caption{밀도별 att\_share 분포 (BRPL vs. TA-BRPL).
  TA-BRPL의 중앙값이 일관되게 낮으며, 중밀도에서 격차가 가장 크다.}
\label{fig:attshare_box}
\end{figure}

\begin{figure}[t]
\centering
\includegraphics[width=0.95\linewidth]{fig_density_bar.pdf}
\caption{밀도별 평균 $\Delta$att\_share(왼쪽)와 $\Delta$PDR(오른쪽) + 95\% CI.
  중밀도에서 두 지표 모두 통계적으로 유의한 개선이 나타난다.}
\label{fig:density_bar}
\end{figure}

\subsection{토폴로지별 격리-안정성 상충 관계}

\begin{figure}[t]
\centering
\includegraphics[width=0.85\linewidth]{fig_scatter.pdf}
\caption{토폴로지별 $\Delta$att\_share vs. $\Delta$churn 산점도.
  왼쪽 아래(격리 개선, churn 비용 소)가 이상적이다.
  대부분의 토폴로지에서 att\_share가 감소하고 churn 증가는 작다.
  저밀도(녹색)에서 일부 큰 churn 사례가 나타난다.}
\label{fig:scatter}
\end{figure}

그림~\ref{fig:scatter}은 75개 토폴로지별 상충 관계를 보여준다.
대부분의 점이 왼쪽 아래 사분면(격리 개선, churn 소)에 위치하며,
churn이 크게 증가한 경우는 대부분 저밀도 네트워크다.

\subsection{공격자 경유율 누적분포}

\begin{figure}[t]
\centering
\includegraphics[width=0.95\linewidth]{fig_cdf.pdf}
\caption{밀도별 att\_share 누적분포.
  TA-BRPL(실선)이 BRPL(점선)보다 일관되게 왼쪽에 위치
  ---평균이 아니라 분포 전체가 이동했음을 보인다.
  중밀도에서 이동이 가장 뚜렷하다.}
\label{fig:cdf}
\end{figure}

그림~\ref{fig:cdf}는 단순 평균이 아니라
분포 전체의 이동을 보여준다.
TA-BRPL의 CDF가 일관되게 왼쪽에 위치하는 것은
특정 토폴로지만 개선된 것이 아니라
전반적인 분포가 이동했음을 의미한다.

\subsection{공격 단계 시계열 증거}

\begin{figure}[t]
\centering
\includegraphics[width=0.95\linewidth]{fig4_route_exposure_timeseries.pdf}
\caption{공격 단계별 시계열(대표 토폴로지 집계).
  수직 점선은 공격 활성 시점($t=350$\,s)을 나타낸다.
  BRPL은 att\_share가 먼저 상승하고 이후 PDR 저하가 뒤따르는 반면,
  TA-BRPL은 초기 route capture 상승폭을 억제해
  후속 PDR 하락도 완화한다.}
\label{fig:stage_timeseries}
\end{figure}

그림~\ref{fig:stage_timeseries}은 본 논문의 핵심 인과를 보여준다.
즉 ``PDR 변화''보다 ``att\_share 변화''가 먼저 나타나며,
TA-BRPL의 이득은 drop 탐지 강화가 아니라
초기 경로 장악 억제에서 시작된다.

\subsection{PDR 비열세 검증 (H2)}

\begin{figure}[t]
\centering
\includegraphics[width=0.88\linewidth]{fig5_pdr_noninferiority.pdf}
\caption{토폴로지별 PDR 비열세 검증 (마진 $\delta = -0.02$\,pp).
  93.3\%의 토폴로지 쌍이 통과.
  고밀도·중밀도는 전부 통과하며, 실패 5건은 모두 저밀도이다.}
\label{fig:noninf}
\end{figure}

H2는 두 수준에서 채택된다:
(1) \emph{집단 수준}: $\Delta\text{PDR}$의 CI 하한 $+0.003 > -0.02$;
(2) \emph{토폴로지 수준}: 93.3\% 통과.
고밀도·중밀도는 전부 비열세를 만족한다.
저밀도 5건의 실패는 경로 재구성 중 일시적 PDR 하락 때문이다.


% ═══════════════════════════════════════════════════════════════════
\section{논의}
\label{sec:discussion}

\subsection{``PDR만 보면 공격의 절반을 놓친다''}

본 논문의 메인 창(350--650초)은 sinkhole 유인과 드롭이 함께 활성화된 구간이다.
그럼에도 att\_share/hit\_ratio가 PDR보다 선행 지표로 작동한다는 점은
시계열(그림~\ref{fig:stage_timeseries})과 쌍체 분포(표~\ref{tab:stats})에서 확인된다.
다만 ``장악-only에서 PDR은 정상인데 att\_share만 먼저 상승한다''는 인과는
stage-split 확장 평가(\S\ref{sec:extended_eval_plan})로 별도 분리 검증한다.

\subsection{RPL과 BRPL이 비슷한 이유: 역압의 한계}

표~\ref{tab:absolute}에서 RPL과 BRPL의 att\_share가 사실상 동일하다.
이는 BRPL의 역압이 큐 정보에 의존하는데,
공격자가 낮은 큐를 광고하면서 무음 드롭을 수행하면
역압이 오히려 공격자를 ``혼잡 없는 좋은 경로''로 홍보하기 때문이다.
TA-BRPL의 $T_{hon}$이 이 구멍을 정확히 메운다.

\subsection{왜 좋아졌는가: 신호와 게이트의 결합 효과}

개선 인과는 ``신호 분해''와 ``전환 안정화''의 결합으로 해석된다.
먼저 $T_{ctrl}$는 저-rank 유인 단계에서 조기 경보를 제공하고,
$T_{hon}$은 허위 큐 광고를 통해 BRPL 역압을 악용하는 경로를 억제한다.
다음으로 admission/dwell/re-entry 게이트는
신뢰 하락을 즉시 대량 전환으로 연결하지 않도록 완충해
불필요한 진동 churn을 줄인다.
이 구조적 효과는 설계 진화 표(표~\ref{tab:versions})에서
v1(게이트 없음) 대비 v13.12(게이트 정렬)로 갈수록
격리 성능과 안정성 균형이 개선되는 경향과 일치한다.

\subsection{Churn 비용의 정직한 해석}

churn 결과를 두 관점에서 정리한다:
\begin{itemize}
  \item \textbf{통계적 관점}: 단측 검정 기준 $p = 1.0$
    (관측 방향이 ``감소'' 가설과 반대이므로).
    양측 기준으로는 소폭 유의한 증가.
  \item \textbf{실무적 관점}: 0.036회/노드/공격구간.
    900초 시뮬레이션의 공격 300초 구간 동안 노드당 평균 0.036회 추가 전환.
    사전 선언 상한 0.1의 36\% 수준이다.
\end{itemize}
churn 비용이 ``없다''고 주장하지 않는다.
다만 설계 목표인 ``격리 이득 대비 관리 가능한 안정성 오버헤드''는
이 실험 범위에서 충족되었다.

\subsection{메커니즘 인과 검증의 다음 단계}

현재 본문은 ``격리 성능 개선''을 통계적으로 입증했지만,
SCI 심사 관점에서 더 강한 인과 증명을 위해
failure mode를 admission과 retention으로 분해한 지표를 추가할 계획이다.
핵심 지표는
\texttt{new\_attacker\_adoption\_rate},
\texttt{attacker\_parent\_retention\_time},
\texttt{escape\_latency\_after\_attack},
\texttt{attacker\_reentry\_count}이며,
이는 ``덜 채택했는가''와 ``더 빨리 탈출했는가''를 분리해 검증한다.

\subsection{확장 평가 계획(현재 미검증 항목)}
\label{sec:extended_eval_plan}

본 논문은 메인 결론(H1/H2/H3)을 BRPL 쌍체 통계로 닫고,
아래 항목은 \emph{확장 평가 계획}으로 분리한다.
즉, 아래 항목은 현 라운드 결론의 근거로 사용하지 않는다.

\medskip
\noindent\textbf{심사 방어용 필수 3항목(Must-have).}
\textbf{(1) Stage-split}, \textbf{(2) Multi-attacker(2--3, cluster 포함)},
\textbf{(3) 현실 무선모델 1종 추가(UDGM + Distance Loss/MRM 계열)}
를 최우선으로 수행한다.

\begin{table}[t]
\centering
\caption{확장 평가 계획(미검증 항목) 요약}
\label{tab:extended_eval_plan}
\small
\begin{tabularx}{\linewidth}{p{2.1cm}X p{2.4cm}X}
\toprule
실험 축 & 설정(예시) & 주요 지표 & 목적 \\
\midrule
Stage-split & 350--500s sinkhole only, 500--650s drop-on & att\_share/PDR 시계열 & ``장악$\rightarrow$파괴'' 인과 분리 검증 \\
Multi-attacker & attacker 1/2/3, sinkhole cluster(근접 2--3) & att/hit/PDR/churn & 협공 시나리오 reviewer 질의 방어 \\
Radio model & UDGM + Distance Loss(MRM/DGRM 중 1종) & 동일 4지표 + 오탐 변화 & ``현실성 부족'' 공격 차단 \\
Overhead/FP & DIO/DAO, 추가 바이트, blacklist 이벤트 & overhead, FP 비율 & 보안 이득 대비 비용 정량화 \\
강건성 스윕 & drop 0/25/50/75/100, attacker 1--3 & att/hit/PDR/churn & 일반화 범위 확인 \\
동급 baseline & RFTrust/PCC-RPL 중 최소 1종 & 동일 4지표 + 비용 & 비교 타당성(공정 비교) 확보 \\
\bottomrule
\end{tabularx}
\normalsize
\end{table}

\subsection{한계와 Threats to Validity}

\begin{itemize}
  \item \textbf{Internal validity}: 본 결론은
    \emph{단일 공격자, UDGM, 50노드, sinkhole+50\% 드롭, v13.12 고정 파라미터}
    조건에 한정된다.
    4-way baseline + drop sweep(0/25/50/75/100\%)은
    파이프라인 사전 검증(\S\ref{sec:precheck4way})을 통과했으나,
    전수 실행이 완료되기 전까지는 최종 우열 결론으로 일반화하지 않는다.
  \item \textbf{External validity}: UDGM은 실제 페이딩·간섭·hidden-node를
    충분히 반영하지 못하며, 현재 결과는 하드웨어 테스트베드로 교차 검증되지 않았다.
  \item \textbf{Construct validity}: 본문 지표(att\_share, hit\_ratio, PDR, churn)는
    ``경로 장악 억제''를 직접 측정하도록 설계되었지만,
    admission/retention/re-entry의 인과 분해는 추가 지표로 보강할 필요가 있다.
  \item \textbf{Attribution bias}: $T_{fwd}$의 귀책 편향
    (정직한 중계 노드가 공격자의 드롭 영향으로 낮은 신뢰를 받는 현상)은
    저밀도 구간에서 잔류한다.
  \item \textbf{Parameter sensitivity}: v13.12 파라미터는 고정 토폴로지에서
    선별된 값이며, 랜덤 토폴로지 분포 전체에 대한 공동 최적화가 남아 있다.
\end{itemize}


% ═══════════════════════════════════════════════════════════════════
\section{결론}
\label{sec:conclusion}

본 연구는 sinkhole 공격을 ``경로 장악 → 가용성 파괴''의 2단계 위협으로
정의하고, 방어의 1차 목표를 PDR이 아닌 경로 장악 억제로 설정한
TA-BRPL을 제안·검증했다.

\textbf{핵심 결론은 세 축으로 정리된다.}
(i) 경로 장악 억제: att\_share/hit\_ratio가 유의하게 감소(H1),
(ii) 가용성 보전: 공격 구간 PDR 비열세를 충족(H2),
(iii) 안정성 상한: churn 증가는 사전 선언 상한 내(H3).
즉, TA-BRPL의 기여는 ``더 민감한 탐지''보다
``신뢰-라우팅 결합의 안정화 의사결정''에 있다.

본 결론은 단일 공격자·UDGM·50노드·v13.12 파라미터의
현재 실험 범위에서 성립한다.
stage-split/overhead/오탐 및 동급 보안 baseline 직접 비교는
\S\ref{sec:extended_eval_plan}의 확장 평가로 남는다.

\textbf{향후 연구 방향은 다음 네 가지로 계획한다.}

\begin{description}
  \item[F1 (비교군 확장)]
    본문의 주 통계 결론은 BRPL 중심 비교에 기반한다.
    동일 파이프라인의 SMTRUST 4-way precheck는
    무결성 점검용으로만 완료되었다(\S\ref{sec:precheck4way}).
    향후에는 표준 RPL(MRHOF/OF0)을 기저선으로 추가해
    ``backpressure 이전 대비'' 개선량을 분리하고,
    동급 보안 RPL 계열(RFTrust~\cite{rftrust2021},
    PCC-RPL~\cite{pccrpl2022})과의 직접 비교로
    TA-BRPL의 상대적 위치를 명확히 할 것이다.

  \item[F2 (강건성 검증)]
    본문 주 실험은 단일 공격자·SINK\_DROP50·UDGM 조건에 한정된다.
    드롭률 변화(0\%, 25\%, 50\%, 75\%, 100\%), 공격자 수(1–3),
    비이상 무선 모델(MRM/DGRM, 간섭·hidden node)에 걸친
    파라메트릭 스윕을 통해 방어 효과의 일반화 범위를 정량화할 예정이다.

  \item[F3 (Ablation 확장)]
    본문 \S\ref{sec:ablation_plan}에서 최소 비교군
    (Full / $T_{fwd}$ only / $(T_{fwd}+T_{ctrl})$)의 pre-main 검증을 완료했다.
    후속으로는 $T_{hon}$ only, gate 제거형(re-entry hold off, recent-switch off),
    soft-only/hard-only까지 확장해
    신뢰 성분 기여와 안정화 게이트 기여를 분리 정량화할 계획이다.

  \item[F4 (재현성 및 오픈 아티팩트)]
    시뮬레이션 코드(Contiki-NG 패치, 토폴로지 생성기, 분석 스크립트),
    원시 로그, 집계 CSV를 공개 저장소(GitHub/Zenodo)에 공개할 예정이다.
    실물 TelosB 테스트베드에서의 재현 실험도 병행해
    시뮬레이션 결과의 현실성을 검증한다.
\end{description}


% ═══════════════════════════════════════════════════════════════════
\section*{데이터 및 코드 가용성}
\label{sec:data}

TA-BRPL 구현(Contiki-NG 패치), 토폴로지 생성기, 실험 스크립트,
원시 로그 및 집계 CSV를 공개할 예정이다.
심사 중에는 익명 저장소 링크로 코드/데이터를 제공하고,
수락 후에는 GitHub + Zenodo(DOI)로 영구 공개한다.
공개 범위는 (i) 시나리오/Makefile/파서 스크립트,
(ii) 메인·precheck·ablation 집계 CSV,
(iii) 원시 \texttt{sim.log} 샘플과 결과 재생성 스크립트를 포함한다.
재현 실행 명령은 저장소 루트 \texttt{README}의
``Reproduce main\_ko results'' 섹션에 고정해 제공한다.

% ═══════════════════════════════════════════════════════════════════
\appendix
\section{실험 상세 설정}
\label{app:setup}

\subsection{소프트웨어 및 플랫폼}

\begin{itemize}
  \item \textbf{시뮬레이터}: Contiki-NG (커밋 기반 수정본),
    Cooja 헤드리스 모드.
  \item \textbf{타깃 플랫폼}: Sky/TelosB.
  \item \textbf{RPL 모드}: Storing mode.
  \item \textbf{MAC}: ContikiMAC (기본 채널 검사).
  \item \textbf{무선 모델}: UDGM, 전송 범위 50m, 간섭 범위 60m.
  \item \textbf{시뮬레이션 시간}: 900초 (각 단계별 일정은 표~\ref{tab:schedule}).
\end{itemize}

\subsection{4-way 사전 검증 스냅샷(부록)}

\begin{table}[htbp]
\centering
\caption{4-way drop-sweep 사전 검증 스냅샷(파이프라인 무결성 점검용).
  본 표는 baseline 우열 결론용이 아니며, 부분 실행 시점의
  스냅샷 예시를 부록에 보존한 것이다.
  값은 \texttt{results/random\_topo\_drop\_sweep\_baseline4/summary/snapshot.md}
  집계 평균이다.}
\label{tab:precheck4way_appendix}
\begin{tabular}{lrrrr}
\toprule
프로토콜 & PDR\_dur & att\_share & hit\_ratio & churn \\
\midrule
RPL     & 0.9845 & 0.0689 & 0.0842 & 14.54 \\
BRPL    & 0.9990 & 0.0667 & 0.0693 & 1.29 \\
SMTRUST & 0.5813 & 0.0652 & 0.1826 & 68.47 \\
TABRPL  & 0.9917 & 0.0806 & 0.0905 & 4.44 \\
\bottomrule
\end{tabular}
\end{table}

\subsection{재현성 고정값(코드 스냅샷)}
\label{app:repro_lock}

\begin{table}[htbp]
\centering
\caption{재현성 식별자(실험 코드 스냅샷)}
\label{tab:repro_hash}
\begin{tabular}{ll}
\toprule
항목 & 값 \\
\midrule
실험 루트 리포지토리 커밋 & \texttt{5cd74bf} \\
Contiki-NG 서브트리 커밋 & \texttt{521724e} \\
\bottomrule
\end{tabular}
\end{table}

\begin{table}[htbp]
\centering
\caption{네트워크/라우팅 핵심 고정 파라미터 (\texttt{project-conf.h})}
\label{tab:repro_params}
\begin{tabular}{ll}
\toprule
항목 & 값 \\
\midrule
RPL DIO Imin (\texttt{RPL\_CONF\_DIO\_INTERVAL\_MIN}) & 12 (4096 ms) \\
RPL DIO Doublings (\texttt{RPL\_CONF\_DIO\_INTERVAL\_DOUBLINGS}) & 8 \\
RPL DIO Redundancy (\texttt{RPL\_CONF\_DIO\_REDUNDANCY}) & 10 \\
송신 주기 (\texttt{SEND\_INTERVAL\_SECONDS}) & 30 s \\
공격 워밍업 (\texttt{ATTACK\_WARMUP\_SECONDS}) & 350 s \\
신뢰 가중치 $(w_f,w_c,w_h)$ & (0.5, 0.3, 0.2) \\
$\lambda_p$ (\texttt{BRPL\_CONF\_TRUST\_LAMBDA\_PENALTY}) & 0.45 \\
상대/절대 전환 마진 & 17\%, 90 (raw) \\
부모 최소 체류시간 & 180 s \\
\bottomrule
\end{tabular}
\end{table}

\begin{table}[htbp]
\centering
\caption{무선/매체 모델 고정값 (Cooja UDGM)}
\label{tab:repro_udgm}
\begin{tabular}{ll}
\toprule
항목 & 값 \\
\midrule
Radio medium & UDGM \\
전송 범위 & 50.0 m \\
간섭 범위 & 60.0 m \\
TX 성공비 (\texttt{success\_ratio\_tx}) & 1.0 \\
RX 성공비 (\texttt{success\_ratio\_rx}) & 1.0 \\
링크 속도 (\texttt{ContikiRadio bitrate}) & 250 kbps \\
\bottomrule
\end{tabular}
\end{table}

또한 본 실험 코드에서는 \texttt{project-conf.h}에서 MAC/RDC를 별도 override하지 않으며,
Sky/Cooja 기본 스택(CSMA + ContikiMAC) 설정을 사용한다.

\subsection{실험 단계 스케줄}

코드 스냅샷 기준으로 공격 노드는 $t=350$\,s 이후 계속 활성화된다
(\texttt{ATTACK\_WARMUP\_SECONDS=350}, 이후 \texttt{attack\_enabled=1} 유지).
분석 파이프라인은 기존 호환성을 위해 650--900초 구간 라벨을
\texttt{recovery}로 유지하지만, 이는 ``공격 오프라인''이 아니라
``공격 후반 관측 창''으로 해석해야 한다.

\begin{table}[htbp]
\centering
\caption{실험 단계 스케줄}
\label{tab:schedule}
\begin{tabular}{lll}
\toprule
구간 & 단계 & 내용 \\
\midrule
0--150초   & 워밍업    & DODAG 구성 \\
150--350초 & 기준선    & 안정 측정 (PDR, att\_share 기준값) \\
200--300초 & 혼잡 유발 & 전송 속도 2배 (큐 동역학 스트레스) \\
350--650초 & 공격 활성(초반) & sinkhole + 드롭($p=\texttt{ATTACK\_DROP\_PCT}$) \\
650--900초 & 공격 활성(후반) & 동일 공격 지속 (분석 라벨: recovery) \\
\bottomrule
\end{tabular}
\end{table}

\subsection{지표 정의}
\label{app:metrics}

지표 집계는 파서의 고정 창(\texttt{pre}: 150--350, \texttt{during}: 350--650,
\texttt{recovery}: 650--900)을 따른다.
이때 \texttt{during}을 본문의 주 공격 평가 창으로 사용한다.

\begin{table}[htbp]
\centering
\caption{평가 지표 정의 및 계산 방법}
\label{tab:metrics}
\begin{tabularx}{\linewidth}{lX}
\toprule
지표 & 정의 및 계산 방법 \\
\midrule
att\_share
  & 공격 활성 구간(350--650초)에 루트가 수신한 패킷 중
    공격자 노드를 경유한 패킷의 비율.
    경로는 시뮬레이션 로그의 홉별 전달 기록에서 추출.
    (\emph{공격자 의존 강도} 지표) \\[4pt]
hit\_ratio
  & 공격 활성 구간에 데이터를 전송한 소스 노드 중
    최소 한 번이라도 공격자를 경유한 패킷을 전송한
    소스 노드의 비율.
    (\emph{공격 노출 범위} 지표) \\[4pt]
PDR\_dur
  & 공격 활성 구간 동안
    (루트 수신 패킷 수) / (소스 전송 패킷 수)의 비율.
    소스별로 먼저 계산 후 평균. \\[4pt]
churn
  & 공격 활성 구간 동안 노드당 선호 부모 변경 횟수의 평균.
    부모 변경은 로그의 부모 갱신 이벤트로 집계. \\
\bottomrule
\end{tabularx}
\end{table}

\subsection{토폴로지 생성}

\begin{itemize}
  \item \textbf{고밀도}: 50노드, $200\times200$m.
  \item \textbf{중밀도}: 50노드, $250\times250$m.
  \item \textbf{저밀도}: 50노드, $300\times300$m.
\end{itemize}
연결성 필터: 루트 도달 가능성, 최대 차수 $\le 15$.
루트는 중앙 고정,
공격자는 betweenness centrality 최고 비루트 노드.

\subsection{최종 파라미터 (v13.12)}

\begin{table}[htbp]
\centering
\caption{TA-BRPL 최종 파라미터 (v13.12)}
\label{tab:params}
\begin{tabular}{ll}
\toprule
파라미터 & 값 \\
\midrule
$\tau_{warn}$ (리뷰 진입) & 600/1000 \\
$\tau_{join}$ (신규 부모 허용) & 510/1000 \\
$\tau_{black}$ (블랙리스트) & 200/1000 \\
$\lambda_{dec}$ (하락 EWMA) & 0.4 \\
$\lambda_{norm}$ (회복 EWMA) & 0.7 \\
신뢰 업데이트 간격 & 60초 \\
블랙리스트 유지 시간 & 120초 \\
전환 마진 (상대) & 17\% \\
전환 마진 (절대) & 90 단위 \\
체류 시간 & 180초 \\
탈출 쿨다운 & 360초 \\
\bottomrule
\end{tabular}
\end{table}

\subsection{통계 방법}

각 토폴로지에서 5 시드의 평균으로 프로토콜별 대표값을 계산하고,
쌍별 차이 $\Delta = \text{TA-BRPL} - \text{BRPL}$를 구한다.
att\_share·hit\_ratio·churn에는 단측 Wilcoxon 부호 순위 검정
(귀무가설: 차이의 중앙값 $= 0$, 대립: 감소 방향),
PDR에는 단측 Wilcoxon (증가 방향)을 적용한다.
4개 지표에 Holm-Bonferroni 다중 비교 보정.
95\% CI는 쌍별 차이의 부트스트랩(10,000회 재표본).
효과크기는 정렬된 순위 이연 상관계수 $r_{rb}$로 보고한다.

\subsection{투고 체크리스트 메모 (Sensors 우선)}
\label{app:submission_memo}

\begin{itemize}
  \item \textbf{원고 정책}: 본문 결론은 BRPL 쌍체 통계(H1/H2/H3)만으로 닫고,
    precheck/확장 평가 결과는 결론 근거에서 분리한다.
  \item \textbf{재현성}: 커밋 해시, Trickle/UDGM/핵심 파라미터, 실행 명령 위치를
    본문+\S\ref{app:repro_lock}+\texttt{README}에 동시 고정한다.
  \item \textbf{데이터 가용성}: 심사 중 익명 링크, 수락 후 공개 저장소 + DOI 부여를 명시한다.
  \item \textbf{도표 정책}: 개념도(Fig.~1)와 통계 그래프를 명확히 구분하고,
    그림 캡션에서 ``개념도/통계근거'' 성격을 혼동 없이 표시한다.
  \item \textbf{필수 3항목(우선순위 1)}:
    (i) stage-split(장악 단계 vs drop 단계 분리),
    (ii) multi-attacker(1$\rightarrow$2--3, sinkhole cluster 포함),
    (iii) radio model 1종 추가(UDGM + Distance Loss/MRM 계열).
  \item \textbf{확장 실험(우선순위 2)}: overhead/FP, RFTrust/PCC-RPL 직접 비교는
    다음 라운드 실험 항목으로 유지한다.
  \item \textbf{대체 저널 대비}: IEEE Access 제출 요소(템플릿 원본+PDF 일치,
    ORCID, 저자 biography, AI 텍스트 공개)는 본문이 아니라
    별도 투고 체크리스트 항목으로만 관리한다.
\end{itemize}


% ═══════════════════════════════════════════════════════════════════
\bibliographystyle{plain}
\bibliography{references}

\end{document}
